%% kapa: exam=1 round=31 year=2020 period=2 form=A subject=appraisal_law subject_ko=감정평가관계법규 seq=13 number=13
%% source: Q-Net 감정평가사 기출문제 (cst00309) artlSeq=5207219 file=extracted/1008/5207219/A형/A형/2020년도 제31회 감정평가사 1차 시험_2교시_A형.pdf
%% conversion: kordoc 4.24.51 -> tools/gen1.py (auto); review=proofread-sonnet; flags=none
\begin{question}{13}
국토의 계획 및 이용에 관한 법령상 도시계획위원회에 관한 설명으로 옳은 것은?
\begin{choices}
  \item 시ㆍ도도시계획위원회는 위원장 및 부위원장 각 1명을 포함한 20명 이상 25명 이하의 위원으로 구성한다.
  \item 시ㆍ도도시계획위원회의 위원장과 부위원장은 위원 중에서 해당 시ㆍ도지사가 임명 또는 위촉한다.
  \item 중앙도시계획위원회의 회의는 재적위원 과반수의 출석으로 개의하고, 출석위원 과반수의 찬성으로 의결한다.
  \item 시ㆍ군ㆍ구도시계획위원회에는 분과위원회를 둘 수 없다.
  \item 중앙도시계획위원회 회의록은 심의 종결 후 3개월 이내에 공개 요청이 있는 경우 원본을 제공하여야 한다.
\end{choices}
\end{question}
\answer{3}
